% !TEX program = xelatex
% English version of schemat_polaczen_DevKit_U10_Q11_v2.tex (pages 1-2) for the Q11 factory. AMPERE POINT, 2026-09-28.
\documentclass[10pt]{article}
\usepackage{fontspec}
\setmainfont{Calibri}
\newfontfamily\sym{Segoe UI Symbol}
\usepackage[table]{xcolor}
\usepackage{geometry}
\geometry{a4paper,landscape,margin=0.8cm}
\usepackage{tikz}
\usetikzlibrary{arrows.meta}
\usepackage{booktabs,array,longtable}
\pagestyle{empty}
\setlength{\parindent}{0pt}

\definecolor{txblue}{HTML}{1D4ED8}
\definecolor{rxgreen}{HTML}{15803D}
\definecolor{pwrred}{HTML}{B91C1C}
\definecolor{ncgrey}{HTML}{6B7280}
\definecolor{copper}{HTML}{C2410C}
\definecolor{padfill}{HTML}{F59E0B}
\definecolor{boardfill}{HTML}{F3F4F6}
\definecolor{modblue}{HTML}{DBEAFE}
\definecolor{x1purple}{HTML}{7C3AED}

\newcommand{\ra}{{\sym →}}
\newcommand{\la}{{\sym ←}}
\newcommand{\approxs}{{\sym ≈}}

\begin{document}

%======================================================================
% PAGE 1: WIRING
%======================================================================
\begin{tikzpicture}[x=1cm,y=1cm]

\node[anchor=north west, text width=27.5cm] at (0,19.3) {%
{\Large\bfseries Wiring diagram: Shelly DevKit M1 \ra\ site U10 (WBR3 footprint) on the Q11 control board}\\[3pt]
{\small Control board X2Q322\_K\_V4 (260703), controller X2Q311OW. Q11 board viewed from the component side, board edge with the module antenna at the bottom.
DevKit viewed from the top, as in the photo. DevKit powered from the Q11 board through a removable jumper.\\ \mbox{AMPERE POINT}, 28 September 2026, EN v1 (from PL v2 of 24 September 2026).}};

% ---------------- DevKit ----------------
\draw[rounded corners=4pt, fill=boardfill, draw=black!70, line width=0.9pt] (2.0,9.6) rectangle (20.4,15.8);
% USB-C
\draw[fill=black!30, draw=black] (1.2,12.3) rectangle (2.3,13.5);
\node[rotate=90, font=\scriptsize, align=center, text=pwrred] at (0.55,12.9) {\textbf{USB-C only with}\\\textbf{the jumper removed}};
% USB-UART bridge
\draw[fill=white, draw=black!60] (3.3,12.7) rectangle (5.7,14.3);
\node[font=\scriptsize, align=center] at (4.5,13.5) {USB--UART bridge\\CP210x};
\draw[black!60] (2.3,12.9) -- (3.3,13.1);
% buttons
\draw[fill=black!35] (2.65,14.05) circle (0.24); \node[font=\tiny] at (2.65,13.6) {BOOT};
\draw[fill=black!35] (2.65,11.35) circle (0.24); \node[font=\tiny] at (2.65,11.8) {RST};
% LED
\draw[fill=white, draw=black!70] (13.7,13.3) circle (0.22);
\node[font=\scriptsize, anchor=west] at (13.95,13.3) {RGB LED};
% MOD1 module
\draw[fill=black!20, draw=black] (16.4,12.6) rectangle (19.4,14.5);
\draw[fill=modblue, draw=black, dashed] (19.4,12.9) rectangle (21.0,14.2);
\node[font=\scriptsize, align=center] at (17.9,13.55) {Shelly-X-MOD1-H8\\ESP32-C3 core};
\node[font=\tiny, rotate=90] at (20.2,13.55) {antenna};
\node[font=\small\bfseries, align=center] at (9.6,13.55) {Shelly DevKit M1\\[-1pt]{\scriptsize\mdseries powered from the Q11 board or from USB, never both}};

% top row (J1), from the left
\foreach \i/\lab/\desc in {
 0/GND/ground, 1/5V/{5 V, unused}, 2/5V/{5 V, unused}, 3/GND/ground, 4/10/{GPIO10, unused},
 5/1/{GPIO1, unused}, 6/0/{GPIO0, unused}, 7/GND/ground, 8/RST/{chip reset}, 9/GND/ground,
 10/3/{GPIO3, unused}, 11/2/{GPIO2, strapping}, 12/3V3/{3.3 V, unused}, 13/3V3/{3.3 V \la\ VCC}, 14/GND/ground}{
  \pgfmathsetmacro{\x}{3.8+1.1*\i}
  \draw[fill=black!45, draw=black] (\x-0.17,15.13) rectangle (\x+0.17,15.47);
  \node[font=\scriptsize\bfseries] at (\x,14.85) {\lab};
  \node[font=\scriptsize, rotate=90, anchor=west, text=black!65] at (\x,16.0) {\desc};
}
\node[font=\scriptsize\itshape, text=black!60, anchor=east] at (3.5,16.6) {top row};

% bottom row (J3), from the left
\foreach \i/\lab/\desc in {
 0/GND/ground, 1/19/{USB D+}, 2/18/{USB D−}, 3/GND/ground, 4/4/{TX \ra\ Q11},
 5/5/{RX \la\ Q11}, 6/6/unused, 7/7/unused, 8/GND/ground, 9/8/{RGB LED},
 10/9/{BOOT}, 11/GND/ground, 12/RX/{log, USB}, 13/TX/{log, USB}, 14/GND/ground}{
  \pgfmathsetmacro{\x}{3.8+1.1*\i}
  \draw[fill=black!45, draw=black] (\x-0.17,9.93) rectangle (\x+0.17,10.27);
  \node[font=\scriptsize\bfseries] at (\x,10.55) {\lab};
  \node[font=\scriptsize, rotate=90, anchor=west, text=black!65] at (\x,10.95) {\desc};
}
\node[font=\scriptsize\itshape, text=black!60, anchor=east] at (3.5,9.35) {bottom row};

% highlighted pins
\draw[fill=txblue, draw=black] (8.2-0.19,9.91) rectangle (8.2+0.19,10.29);
\draw[fill=rxgreen, draw=black] (9.3-0.19,9.91) rectangle (9.3+0.19,10.29);
\draw[fill=black, draw=black] (12.6-0.19,9.91) rectangle (12.6+0.19,10.29);
\draw[fill=pwrred, draw=black] (18.1-0.19,15.11) rectangle (18.1+0.19,15.49);
% RX/TX on the header are the USB/log port
\foreach \x in {17.0,18.1}{
  \draw[pwrred, line width=1.6pt] (\x-0.28,9.82) -- (\x+0.28,10.38);
  \draw[pwrred, line width=1.6pt] (\x-0.28,10.38) -- (\x+0.28,9.82);
}
\node[font=\scriptsize\bfseries, text=pwrred, align=center, fill=white, inner sep=1.5pt] at (17.55,9.3) {RX and TX: NOT THESE!\\USB and log port};
\node[font=\scriptsize\bfseries, text=txblue, fill=white, inner sep=1.5pt] at (8.75,9.3) {wires to the Q11: pins 4 and 5};

% ---------------- site U10 ----------------
\fill[modblue] (14.6,0.0) rectangle (20.6,0.7);
\node[font=\scriptsize, text=black!70] at (17.6,0.35) {module antenna over the board cut-out, board edge};
\draw[fill=white, draw=black, line width=0.9pt] (14.6,0.7) rectangle (20.6,8.9);
\node[align=center, font=\small] at (17.6,4.8) {\textbf{U10}\\WBR3\\footprint};
\fill (21.05,0.75) circle (0.07);
\node[font=\tiny, anchor=west] at (21.15,0.75) {dot on the board = pad 1};

% left column: 9..16 from the top
\foreach \k/\n/\name in {0/9/GND,1/10/A\_12,2/11/A\_16,3/12/A\_17,4/13/A\_18,5/14/A\_19,6/15/RXD,7/16/TXD}{
  \pgfmathsetmacro{\y}{8.3-1.0*\k}
  \draw[fill=padfill, draw=copper] (14.25,\y-0.18) rectangle (14.95,\y+0.18);
  \node[anchor=west, font=\scriptsize\bfseries] at (15.05,\y) {\n\ \ \name};
}
% right column: 8..1 from the top
\foreach \k/\n/\name in {0/8/VCC,1/7/A\_4,2/6/A\_3,3/5/A\_2,4/4/A\_11,5/3/EN,6/2/A\_7,7/1/NC}{
  \pgfmathsetmacro{\y}{8.3-1.0*\k}
  \draw[fill=padfill, draw=copper] (20.25,\y-0.18) rectangle (20.95,\y+0.18);
  \node[anchor=east, font=\scriptsize\bfseries] at (20.15,\y) {\name\ \ \n};
}

% unconnected pads, left column (10..14)
\foreach \k in {1,2,3,4,5}{
  \pgfmathsetmacro{\y}{8.3-1.0*\k}
  \draw[ncgrey, line width=0.8pt] (14.25,\y) -- (13.85,\y);
  \node[font=\scriptsize, text=ncgrey, anchor=east] at (13.8,\y) {× not connected};
}

% right column descriptions (outside)
\draw[pwrred, line width=1.7pt] (20.95,8.3) -- (21.6,8.3) -- (21.6,11.25);
\draw[pwrred, line width=1.7pt] (21.6,12.15) -- (21.6,15.65) -- (18.1,15.65) -- (18.1,15.49);
\draw[pwrred, line width=1.2pt, fill=white] (21.6,11.35) circle (0.1);
\draw[pwrred, line width=1.2pt, fill=white] (21.6,12.05) circle (0.1);
\draw[pwrred, line width=1.2pt, fill=white] (21.35,11.25) rectangle (21.85,12.15);
\fill[pwrred] (21.6,11.35) circle (0.06);
\fill[pwrred] (21.6,12.05) circle (0.06);
\node[anchor=west, font=\scriptsize, text=pwrred, text width=5.2cm] at (21.95,11.7) {\textbf{jumper}: remove before connecting USB; insert and remove only with the charger switched off};
\node[anchor=west, font=\scriptsize\bfseries, text=pwrred] at (21.7,9.2) {3.3 V from the board \ra\ DevKit 3V3};
\node[anchor=west, font=\scriptsize\bfseries, text=pwrred] at (21.7,14.6) {to the top row, pin 14: 3V3};
\foreach \k/\t in {1/{module GPIO}, 2/{module GPIO}, 3/{module GPIO}, 4/{module GPIO}, 5/{WBR3 enable; the DevKit has its own RST}, 6/{module GPIO}}{
  \pgfmathsetmacro{\y}{8.3-1.0*\k}
  \draw[ncgrey, line width=0.8pt] (20.95,\y) -- (21.3,\y);
  \node[anchor=west, font=\scriptsize, text=ncgrey] at (21.35,\y) {× not connected: \t};
}
\node[anchor=west, font=\scriptsize, text=ncgrey] at (21.1,1.3) {× pad not connected (NC)};

% ---------------- wires ----------------
\draw[txblue, line width=1.7pt] (8.2,9.91) -- (8.2,1.3) -- (14.25,1.3);
\draw[rxgreen, line width=1.7pt] (9.3,9.91) -- (9.3,2.3) -- (14.25,2.3);
\draw[black, line width=1.7pt] (12.6,9.91) -- (12.6,8.3) -- (14.25,8.3);
\draw[fill=white, draw=txblue, line width=1.2pt] (8.02,7.35) rectangle (8.38,8.35);
\draw[fill=white, draw=rxgreen, line width=1.2pt] (9.12,7.35) rectangle (9.48,8.35);
\node[font=\scriptsize\bfseries, text=txblue, rotate=90] at (7.8,7.85) {470 {\sym Ω}};
\node[font=\scriptsize\bfseries, text=rxgreen, rotate=90] at (9.7,7.85) {470 {\sym Ω}};
\node[font=\tiny, text=black!60, rotate=90] at (9.98,7.85) {recommended, note 3};
% direction arrows
\draw[-{Stealth[length=3.4mm,width=2.6mm]}, txblue, line width=1.7pt] (12.2,1.3) -- (13.4,1.3);
\draw[-{Stealth[length=3.4mm,width=2.6mm]}, rxgreen, line width=1.7pt] (9.3,5.2) -- (9.3,6.5);
\draw[-{Stealth[length=3.4mm,width=2.6mm]}, txblue, line width=1.7pt] (8.2,6.5) -- (8.2,5.2);
% wire labels
\node[font=\scriptsize\bfseries, text=txblue, anchor=south west] at (8.35,1.33) {DevKit pin 4, TX output \ra\ pad TXD \ra\ controller RX input (PA3)};
\node[font=\scriptsize\bfseries, text=rxgreen, anchor=south west] at (9.45,2.33) {DevKit pin 5, RX input \la\ pad RXD \la\ controller TX output (PA2)};
\node[font=\scriptsize\bfseries, anchor=south] at (13.45,8.33) {ground};

% controller: direction note
\node[draw=black!50, dashed, rounded corners=3pt, align=center, font=\scriptsize, text width=2.3cm, fill=white] at (11.0,5.6) {%
\textbf{Q11 controller}\\(X2Q311OW) on the same board, connected by tracks:\\[2pt]
pad TXD \ra\ its RX\\pad RXD \la\ its TX};

% ---------------- legend and notes ----------------
\node[anchor=north west, draw=black!40, rounded corners=3pt, fill=white, text width=6.75cm, font=\scriptsize, inner sep=5pt] at (0.2,9.15) {%
\textbf{Legend}\\[2pt]
{\color{txblue}\rule[0.5ex]{0.8cm}{1.7pt}}\ transmit: DevKit \ra\ controller\\
{\color{rxgreen}\rule[0.5ex]{0.8cm}{1.7pt}}\ receive: controller \ra\ DevKit\\
{\color{black}\rule[0.5ex]{0.8cm}{1.7pt}}\ common ground\\
\fbox{\rule{0pt}{0.8ex}\hspace{0.6cm}}\ 470 {\sym Ω} series resistor, at the DevKit\\
{\color{pwrred}\rule[0.5ex]{0.8cm}{1.7pt}}\ 3.3 V from the board, through the jumper\\
{\color{ncgrey}×}\ pad left open\\[4pt]
\textbf{Notes}\\[2pt]
1. Solder the wires to the DevKit pins labelled \textbf{4 and 5}, the fifth and sixth from the USB-C connector in the bottom row. \textbf{Not to the pins labelled RX and TX} at the end of that row: they are the DevKit's USB and log port. Pin 4 to pad TXD, pin 5 to pad RXD, no crossing.\\[2pt]
2. \textbf{Power from the board or from USB, never both.} Insert and remove the jumper only with the charger switched off.\\[2pt]
3. The 470 {\sym Ω} resistors are recommended; they protect against a wiring mistake (3.3 V / 470 {\sym Ω} \approxs\ 7 mA) and do not affect the signal (470 {\sym Ω} × 20 pF \approxs\ 10 ns against a 104 µs bit at 9600 baud).\\[2pt]
4. Before soldering, charger on: VCC–GND \approxs\ 3.3 V, RXD–GND \approxs\ 3.3 V (idle high).\\[2pt]
5. In operation, with Wi-Fi active: VCC–GND at least 3.2 V.\\[2pt]
6. In the Shelly portal: IO4 = TX, IO5 = RX; IO number = pin label.\\[2pt]
7. Pads A\_x and EN: measured, no continuity to anything on this board.\\[2pt]
8. DevKit inside the enclosure: insulated, fixed, away from mains-voltage parts.};

\end{tikzpicture}

\newpage
%======================================================================
% PAGE 2: SITE U10 - FUNCTION OF EVERY PAD
%======================================================================
\begin{tikzpicture}[x=1cm,y=1cm]

\node[anchor=north west, text width=27.5cm] at (0,19.3) {%
{\Large\bfseries Site U10: function of every pad with the WBR3 \ra\ with the DevKit \ra\ with the Shelly X1 at the same position}\\[3pt]
{\small Component side, board edge with the antenna at the bottom. Left column: pads 9–16 from the top. Right column: pads 1–8 from the bottom, dot next to ``U10'' = pad 1.
Pad names from the WBR3 datasheet. X1 pad names from the Shelly X1 datasheet v14 (IO layout), derived by overlaying the two footprints; to be verified on a physical X1 sample.}};

\fill[modblue] (10.7,0.6) rectangle (16.9,1.6);
\node[font=\footnotesize, text=black!70] at (13.8,1.1) {module antenna over the board cut-out};
\draw[fill=white, draw=black, line width=1pt] (10.7,1.6) rectangle (16.9,17.2);
\node[align=center, font=\large] at (13.8,9.4) {\textbf{U10}\\[2pt]\normalsize module site\\\normalsize 2 × 8 pads, 2 mm pitch};
\fill (17.5,1.8) circle (0.09);
\node[font=\scriptsize, anchor=west] at (17.65,1.8) {dot = pad 1};

% left column
\foreach \k/\n/\name in {0/9/GND,1/10/A\_12,2/11/A\_16,3/12/A\_17,4/13/A\_18,5/14/A\_19,6/15/RXD,7/16/TXD}{
  \pgfmathsetmacro{\y}{16.4-2.0*\k}
  \draw[fill=padfill, draw=copper] (10.25,\y-0.25) rectangle (11.15,\y+0.25);
  \node[anchor=west, font=\normalsize\bfseries] at (11.3,\y) {\n\ \ \name};
}
% right column
\foreach \k/\n/\name in {0/8/VCC,1/7/A\_4,2/6/A\_3,3/5/A\_2,4/4/A\_11,5/3/EN,6/2/A\_7,7/1/NC}{
  \pgfmathsetmacro{\y}{16.4-2.0*\k}
  \draw[fill=padfill, draw=copper] (16.45,\y-0.25) rectangle (17.35,\y+0.25);
  \node[anchor=east, font=\normalsize\bfseries] at (16.3,\y) {\name\ \ \n};
}

% left column descriptions
\newcommand{\gpioL}[1]{{\color{black!55}WBR3: module GPIO. No continuity to anything on the Q11 board.}\\\ra\ {\color{ncgrey}\textbf{DevKit: not connected.}}\\\ra\ {\color{x1purple}\textbf{X1 here: #1}}}
\foreach \k/\t in {
 0/{{\color{black!55}WBR3: module ground.}\\\ra\ \textbf{DevKit: ground. Wire to any DevKit GND, e.g. bottom row between 7 and 8.}\\\ra\ {\color{x1purple}\textbf{X1 here: GND}}},
 1/{\gpioL{Sys btn, left open}}, 2/{\gpioL{U0TXD, service UART, left open}}, 3/{\gpioL{U0RXD, service UART, left open}}, 4/{\gpioL{IO7, unused}}, 5/{\gpioL{IO6, unused}},
 6/{{\color{black!55}WBR3: module serial input. Receives Tuya frames from the controller TX output (PA2).}\\\ra\ {\color{rxgreen}\textbf{DevKit: pin 5 (IO5, RX), 3.3 V input, via 470 {\sym Ω}. Receives the same frames.}}\\\ra\ {\color{x1purple}\textbf{X1 here: IO5 = RX}}},
 7/{{\color{black!55}WBR3: module serial output. Sends Tuya frames to the controller RX input (PA3).}\\\ra\ {\color{txblue}\textbf{DevKit: pin 4 (IO4, TX), 3.3 V output, via 470 {\sym Ω}. Sends frames instead of the WBR3.}}\\\ra\ {\color{x1purple}\textbf{X1 here: IO4 = TX}}}}{
  \pgfmathsetmacro{\y}{16.4-2.0*\k}
  \draw[black!30] (10.25,\y) -- (10.05,\y);
  \node[anchor=east, align=right, text width=9.6cm, font=\footnotesize] at (10.0,\y) {\t};
}

% right column descriptions
\newcommand{\gpioR}[1]{{\color{black!55}WBR3: module GPIO. No continuity to anything on the Q11 board.}\\\ra\ {\color{ncgrey}\textbf{DevKit: not connected.}}\\\ra\ {\color{x1purple}\textbf{X1 here: #1}}}
\foreach \k/\t in {
 0/{{\color{black!55}WBR3: 3.3 V supply input from the control board regulator; reaches controller pins 1, 13, 19, 32, 48, 64.}\\\ra\ {\color{pwrred}\textbf{DevKit: through the jumper to DevKit 3V3, top row, pin 14.}} Powers the DevKit instead of USB.\\\ra\ {\color{x1purple}\textbf{X1 here: 3V3}}},
 1/{\gpioR{IO3, unused}}, 2/{\gpioR{IO2, unused}}, 3/{\gpioR{IO1, unused}}, 4/{\gpioR{IO0, unused}},
 5/{{\color{black!55}WBR3: module enable input (high = run, low = module off). No continuity to anything on the Q11 board: the controller cannot switch the module off.}\\\ra\ {\color{ncgrey}\textbf{DevKit: not connected.} The DevKit has its own reset (RST).}\\\ra\ {\color{x1purple}\textbf{X1 here: NC}}},
 6/{\gpioR{NTC, left open}},
 7/{{\color{black!55}WBR3: not connected (NC).}\\\ra\ {\color{ncgrey}\textbf{DevKit: nothing.}}\\\ra\ {\color{x1purple}\textbf{X1 here: RESET, left open}}}}{
  \pgfmathsetmacro{\y}{16.4-2.0*\k}
  \draw[black!30] (17.35,\y) -- (17.55,\y);
  \node[anchor=west, align=left, text width=9.8cm, font=\footnotesize] at (17.6,\y) {\t};
}

\node[anchor=south west, font=\footnotesize] at (0.2,0.2) {%
Colours: {\color{txblue}\textbf{blue}} transmit DevKit \ra\ controller,
{\color{rxgreen}\textbf{green}} receive controller \ra\ DevKit, \textbf{black} ground,
{\color{pwrred}\textbf{red}} 3.3 V from the board to the DevKit, {\color{ncgrey}\textbf{grey}} pad left open, {\color{x1purple}\textbf{purple}} Shelly X1 pad at the same position.};

\end{tikzpicture}

\end{document}
